\subsection{Type 11 and the three disks at 37}\label{app:type11-details}

We give the global squareclass argument and the Chabauty calculation
underlying Appendix~\ref{app:type11}. Retain its center-square variables
$a,b,c$ and the sum and product $S,T$ of the two remaining squared
centers. In the two higher moment equations the linear subresultant
in $S$, modulo the conic of that appendix, is a nonzero multiple of
\[
 \bigl(16bc+4b-16c^2+56c-85\bigr)
                 \left(S-\frac{5b-4c+16}{3}\right).
\]
The resultant of the first factor with the conic in $b$ is
$16(4c-49)(4c-9)^2(4c-1)$. The exceptional values $1/4,9/4$
meet the central chain. At $c=49/4$ the conic gives $b=9,-123/2$.
The only nonnegative square choice $b=9$ gives $S=4,5$ in the two
uncanceled equations, respectively forcing $a=1/4,9/4$. These also
repeat a central vertex. This proves the displayed formula for $S$
for every valid input and justifies the slope through $(b,c)=(9,49/4)$.

For $w=(b-9)/(c-49/4)$ put
\[
\begin{aligned}
 D&=w^2+5w-5,& F_A&=25w^2-53w+27,\\
 F_B&=41w^2-70w+30,&F_C&=49w^2-37w-5,\\
 F_E&=(4w-5)(47w-40).
\end{aligned}
\]
The necessary center data are
\begin{equation}\label{eq:md11-param}
 a=-\frac{15F_A}{4D},\quad b=-\frac{3F_B}{2D},\quad
 c=\frac{F_C}{4D},\quad T=\frac{F_E(5w-4)^2}{2D^2}.
\end{equation}
The quadratic $F_B$ is strictly positive on the real line, so $b>0$
forces $D<0$. Write $w=m/n$ primitively with $n>0$, and homogenize
all quadratic forms to the pair $(m,n)$. A prime $p\notin\{2,3,5\}$
dividing $D_h$ does not divide $n$. Since
$\operatorname{Res}(D,F_A)=-36$, it does not divide $F_{A,h}$,
so $v_p(D_h)$ is even by the square condition on $a$. Primitivity
makes $D_h$ odd. If $3\mid D_h$, then $3\nmid n$ and
\[
 4D_h=9\left(\left(\frac{2m+5n}{3}\right)^2-5n^2\right).
\]
The parenthesis is a 3-adic unit because 5 is nonsquare modulo 3.
Consequently $v_3(D_h)=2$. The negative squarefree part of $D_h$
is therefore $d=-1$ or $-5$.

The four integer forms
$dD_h,-15dF_{A,h},-6dF_{B,h},dF_{C,h}$ are squares.
The complete projective test modulo 37 leaves slopes $23,30$ for
$d=-1$ and slopes $0,4$ for $d=-5$. The condition $T\ne0$
requires $2F_{E,h}$ to be square. At the first two slopes its squareclass
is represented by 19 and 8, both nonsquares modulo 37. Hence
\begin{equation}\label{eq:md11-class}
 D_h=-5k^2,\qquad w\equiv0,4\pmod{37}.
\end{equation}
The projective test also excludes $n\equiv0\pmod{37}$, which proves
integrality of $w$ at the prime used below.

Set $r=(47w-40)/(4w-5)$. The denominator cannot vanish for a valid
input because it would make $T=0$. The squareclasses in
\eqref{eq:md11-param} give a point on
\[
 C:Y^2=f(r)=r(r^2-11r-1)(7r^2-2r+18)(9r^2-69r-49).
\]
The residual slopes in \eqref{eq:md11-class} require $r\equiv8,0$
at 37. The curve has good reduction there, so precisely three disks
are eligible, containing $(0,0),(8,1500),(8,-1500)$.

A full two-Selmer computation with Minkowski proof bounds gives
dimension five. The four irreducible factors of degrees $1,2,2,2$
contribute three independent rational two-torsion classes. Thus
$\operatorname{rank}J(C)(\mathbf Q)\leq2$.
Define
\[
 P=[(8,1500)-\infty],\qquad
 D=\left(r^3-\frac{51}{4}r^2+\frac{23}{4}r+8,
              \frac{415}{8}r^2+\frac{785}{8}r+20\right).
\]
The second expression is a Mumford divisor with $Y=v(r)$ on its
support. Its support polynomial is squarefree, coprime to $f$ and
divides $f-v^2$. For $\omega_j=r^jdr/(2Y)$, $j=0,1,2$, exact
arithmetic modulo $37^2$ gives
\[
 \frac{\log(58624P)}{37}=(25,28,33),\qquad
 \frac{\log(58624D)}{37}=(3,30,12)\pmod{37}.
\]
Their cross product is $(12,21,0)$. Both global logarithms are
independent, so their span is the full rational logarithm space.
There is an annihilating differential $\sum\lambda_j\omega_j$ with
$(\lambda_0,\lambda_1,\lambda_2)\equiv(12,21,0)\pmod{37}$.

To specify the local calculation, translate each multiple by a
nonspecial effective divisor with three distinct abscissas and unit
ordinates. The following table lists its residue coordinates and the
Mumford endpoint of the translated multiple. Polynomial coefficients
are in descending order modulo $1369$.
\[
\begin{array}{c|c|c|c}
R&\text{auxiliary }(r,Y)&u_{B+58624R}&v_{B+58624R}\\ \hline
P&(1,-7),(3,8),(6,3)&(1,693,101,1129)&(1324,576,276)\\
D&(1,7),(6,-3),(8,-17)&(1,540,876,1173)&(808,1113,417)
\end{array}.
\]
Hensel lifting at each fixed abscissa supplies the auxiliary ordinate.
Unit-pivot Cantor composition and reduction~\cite{Cantor} certify the endpoints.
For doubling, the ordinate lift is $v+uh$ with
$h\equiv((f-v^2)/u)(2v)^{-1}\pmod u$. The endpoint support roots
satisfy $r_i'-r_i\equiv-u(r_i)/u'(r_i)\pmod{37^2}$ and lie on
the same sheets. Therefore the tiny integrals are
$(r_i'-r_i)r_i^j/(2Y_i)\pmod{37^2}$. All nonlinear primitive
terms have valuation at least two, which proves the stated logarithms.

On the two disks over $r=8$, the numerator of the annihilator is
$12+21\cdot8=32\pmod{37}$, a unit. On the Weierstrass disk use $Y$
as parameter, so $dr/(2Y)=dY/f'(r)$ and the constant is
$12/f'(0)=12/31\pmod{37}$, also a unit. Each primitive is injective
on its whole disk. Indeed a nonlinear term of degree $k\geq2$
in the difference at parameters $t_1,t_2\in37\mathbf Z_{37}$ has
valuation at least
\[
 v_{37}(t_1-t_2)+k-1-v_{37}(k)>v_{37}(t_1-t_2).
\]
The unit linear term cannot cancel. Integrals between rational points
in one disk vanish, so the three anchors exhaust the eligible disks.
At $r=0$, $w=40/47$ gives $F_E=0$ and $T=0$. At $r=8$, $w=0$
gives $c=1/4$. All three are invalid. The computations are
\cert{check_type11_genus3_rank_bounds.m} and
\cert{check_type11_thirty_seven_adic.py}, with the preceding elimination
and squareclass tables in \cert{check_all_shift_chain_types.py}.

\subsection{Type 13 and six forced Prym directions}\label{app:type13-details}

We expand the global rank contradiction used in
Appendix~\ref{app:type13}. Put
\[
 N(c)=-240c^3+1600c^2-260c+475,\quad
 \alpha=25/4,\quad\beta=1/4,
\]
and let $K=\mathbf Q(\theta)$ with $N(\theta)=0$. The product of
the three first-side singleton squares is $N(c)/144$. The two other
center squares satisfy $b=(25-4c)/6$, $c=w^2$, so every input gives
a point on the genus-two curve
\[
 H:y^2=6c(25-4c)N(c).
\]
Its two-Selmer group bounds the rank by two. The four classes
\[
 P=[(\beta,135)-\infty],\quad
 Q=[(19/4,2565/2)-\infty],\quad
 T_0=[(0,0)-\infty],\quad T_\alpha=[(\alpha,0)-\infty]
\]
are independent modulo doubling. To check this without assuming a
point-search basis, pass to the monic model by $X=5760c$,
$Y=5760^2y$. The two rational-root components of the Kummer images
of these four classes are
\[
 (10,-15),\quad(190,-15),\quad(285,-10),\quad(10,-15).
\]
Valuations at 2,3,19 leave only $P+T_\alpha$ as a possible nonzero
relation. In the cubic component its image is
$(\beta-\theta)/(\alpha-\theta)$. At the simple residue root
$\theta=6\pmod{11}$ the ratio is the nonsquare 10. Thus the four
classes are independent and form a basis modulo doubling. A nonzero
square value of $c$ forces the first rational Kummer coordinate to
have class 10. The same valuation test reduces its class to
$P$ or $T_\alpha$, giving precisely
\begin{equation}\label{eq:md13-twists}
 (c-\theta)/\delta\in K^{\times2},\qquad
          \delta\in\{\alpha-\theta,\beta-\theta\}.
\end{equation}
Their second rational coordinate is $-15$, and hence the condition
that $b$ be square is preserved. The group bounds and Kummer character
calculations are \cert{check_type13_product_covers.m} and
\cert{check_type13_product_covers.py}.

\subsubsection{The alpha class and divisibility at the ramified dyadic place}

The alpha cover is an intersection of quadrics over $K$,
$w^2-\delta z^2=\theta$, $4w^2+6v^2=25$. Taking the kernel over
$d\theta$, with $d=-6\delta$, on its associated elliptic curve gives
\[
 E':Y^2=X(X^2+aX+e^2),\quad
 a=2d(\alpha-2\theta),\quad e=d\alpha,
\]
with rational center map
\[
 c=\frac{\alpha(X^2+aX+e^2)}{(X-e)^2},\qquad
 \frac{\alpha}{c}=1+24\delta^2t^2,\quad t=-X/Y,
 \qquad \delta t=\pm\frac{v}{4w}.
\]
A valid lift makes $XY(X-e)\ne0$.
Use $K=\mathbf Q(u)$ with $u^3+u^2+315u+1983=0$ and
$\theta=(u^2-3u+342)/60$. Its unique dyadic place has uniformizer
$\pi=(u+1)^2/2$ satisfying
$\pi^3+314\pi^2+26632\pi-347778=0$. Normalize $v_\pi(2)=3$.
The valuations in $4w^2+6v^2=25$ give
$v_2(w)=-1$, $m=v_2(v)\geq1$. Also
$v_2(\delta)=-2$ and $v_2(e)=-3$. Put $s=X/e$.
Since $c-\alpha=4\delta s/(s-1)^2$ has valuation $2m-1>0$,
$s$ is not a unit. Translation by $(0,0)$ replaces $s$ by $1/s$
without changing $c$, so arrange $v_2(s)=-(2m-1)$.

On the integral model $x'=\pi^{10}X$, $y'=\pi^{15}Y$, the
coefficients of $x'^2,x'$ have valuations 4 and 2. The point is
formal, with
\[
 t'=-x'/y'=\kappa2^{m-1}h,\qquad
 \kappa=(\delta\pi^5)^{-1},\quad v_\pi(\kappa)=1,
 \quad h\in\mathbf Z_2^\times.
\]
The formal logarithm identifies this formal group $G$ with
$\pi\mathcal O$, and $\log(t')\equiv t'\pmod{4t'\mathcal O}$.
For a direct bound, set $-1/y'=t'^3U$. Then
$U=1+a't'^2U+b't'^4U^2$, and the invariant differential is
$(1+t'U'/(2U))dt'$. Its nonconstant terms lie in the ideal generated
by $a't'^2,b't'^4$, both of valuation at least six. Their powers of
$t'$ are even, so integration divides only by odd integers. The
logarithm also has an integral compositional inverse.

The global two-Selmer group of $E'$ has order four and injects on
localizing at this dyadic place. Its two local torsion points are
global: the other two would require a square root of
$\theta(\theta-\alpha)$, whose norm is
$35^2\cdot19/2^{10}$, a dyadic nonsquare. At the reference parameter
$t'=\kappa$, there is a point $P_*$ with trivial local Kummer
class, but every local half has Kummer class outside the localized
global Selmer image. To specify this assertion, choose
$z_*\equiv1\pmod{\pi^6}$ with $z_*^2=1-6/\delta$, and put
\[
 X_*=(5\delta(1+z_*)/2)^2,\qquad Y_*=-\delta X_*.
\]
It has $\delta(-X_*/Y_*)=1$. Its halves are calculated by the
identity, on $y^2=x(x^2+a_2x+a_4)$,
\[
 r^2=x,\quad p=-y/r,\quad s^2=2x+a_2+2p,\quad
 x_h=x+rs+p,\quad y_h=-y-(r+s)(rs+p).
\]
The square roots are certified by Hensel margins. On the integral
model $x=60^2X$, $y=60^3Y$, coordinate accuracy exceeding
$2v_\pi(y)+6$ preserves all Kummer squareclasses, since the three
factor valuations are nonnegative and sum to $2v_\pi(y)$.
The resulting injectivity and the two reference classes are checked
by \cert{check_type13_alpha_selmer.m}.

It follows that every global point locally divisible by two is
globally divisible by two. Every chosen local half of such a point
is global because all local two-torsion is global. Divide the
putative formal point $m-1$ times within $G$, obtaining a global
point $R$ whose logarithm is $\kappa h\pmod{4\pi\mathcal O}$.
Choose an odd integer $j\equiv h\pmod4$. Then $R-jP_*\in4G$,
so $R$ is locally divisible by two and has a global half. Its local
half class is that of $j$ times a half of $P_*$, up to rational
two-torsion. Both possible classes lie outside the global Selmer
image, a contradiction. This excludes the alpha class on its whole
rational parameter range.

\subsubsection{The signed-sum quartic and its divisor map}

For the beta class set $\delta=\beta-\theta$ and write the singleton
centers as $a,b,d$. Their squared sum $t=(a+b+d)^2$ satisfies
\[
 ((t-e_1)^2-4e_2)^2=64te_3,
\]
where $e_i$ are the elementary symmetric functions of $a^2,b^2,d^2$.
The homogeneous quadratic in $(c,t,Z)$ is
\[
 q_2=t^2-\frac{10c+50Z}{3}t+5c^2-30cZ+19Z^2.
\]
Let $q(c,Z)$ be the degree-two homogenization of
$(-N(c)/240)/(c-\theta)$ and put
$q_1=t(c-\theta Z)$, $q_3=-(320/3)q(c,Z)$.
The smooth plane quartic is $q_2^2=q_1q_3$. The two lines defining
$q_1$ cut twice distinct degree-two divisors. Equality of those
classes would supply a degree-two pencil on a smooth plane quartic,
which is impossible. Adjoining $\sqrt{q_1/(\delta Z^2)}$ therefore
gives a nontrivial étale double cover. Every beta lift maps to it by
\eqref{eq:md13-twists}. If $M_i$ is the symmetric matrix of $q_i$,
its Prym is the Jacobian $J/K$ of
\begin{equation}\label{eq:md13-prym}
 Y^2=-\delta\det(M_1+2XM_2+X^2M_3).
\end{equation}
This is Bruin's determinant construction~\cite[Theorem~5.1]{BruinPrym}, with twist obtained from
$q_1/\delta,q_2,\delta q_3$ and the change $X\mapsto\delta X$.
The twist and smoothness are checked in
\cert{check_type13_product_covers.py}.

Write $H(c,t)=q_2(c,t,1)^2-4tN(c)/9$ and $h=c-\beta$.
The necessary local reduction at 23 forces $h\in23\mathbf Z_{23}$.
For $h=0$ the input collides. The three distinct nonzero singleton
squares at that point allow their rational roots to be labeled and
signed so that $(a,b,d)\equiv(1/2,3/2,5/2)\pmod{23}$.
Use the ordered signed sums
\[
 v_1=-a-b+d,\quad v_2=a-b+d,\quad
 v_3=-a+b+d,\quad v_4=a+b+d.
\]
For every sum take $t=v_i^2$, $r=v_i z$, $s=q_2/r$, with the
same global square root $z^2=(c-\theta)/\delta$. At $h=0,z=1$
the sums are $1/2,3/2,7/2,9/2$. Let $D_i$ be their four Prym
images at this collision, and $A_i$ the images of a putative input.
We shall show that $D_1,D_2,D_3,A_1,A_2,A_3$ are independent
modulo torsion whenever $h\ne0$.

Here are formulas fixing both the divisor images and all signs.
Primes mean differentiation along $H(c,t)=0$. Put
\[
 a_0=t'/\delta-(r')^2,\quad
 b_0=(t')^2-10t'/3+5-r's',\quad
 k_0=-320\delta/3-(s')^2.
\]
The monic Mumford support on \eqref{eq:md13-prym} is
\[
 U=X^2+2\delta b_0X/k_0+\delta^2a_0/k_0.
\]
Its ordinate is the remainder modulo $U$ of
\begin{multline}\label{eq:md13-ordinate}
 \delta^2\left[
 (r'+Xs'/\delta)\left(\frac{c-\theta}{2\delta}
                      -\frac X\delta\left(\frac{10c+50}{3}-2t\right)\right)\right.\\
 \left. +(r+Xs/\delta)\left(-\frac1{2\delta}
                         -\frac X\delta(2t'-10/3)\right)\right].
\end{multline}
To derive it, take $V=(c,t,1)^t$, $W=(1,t',0)^t$ in the pencil
$Q=M_1/\delta-2xM_2+\delta x^2M_3$. Its tangent condition is
$a_0-2b_0x+k_0x^2=0$. Put $q_v=r-xs$, $q_w=r'-xs'$ and
$I=q_wV-q_vW$. At a tangent root, $QI=\lambda(V\times W)$.
The adjugate identity, together with the vector triple product, gives
$\lambda^2=-\det Q$. The second coordinate of $V\times W$ is
one, so it fixes the signed ruling coordinate. Changing to
$X=-\delta x$ and scaling this coordinate by $\delta^2$ gives
\eqref{eq:md13-ordinate}. This is one rational divisor map over $K$.
Changing $z$ to $-z$ negates the image.

\subsubsection{Equal signs and the analytic determinant}

Order the three split places over 23 by
$\theta\equiv11,13,21$. The reduction orders of $J$ are
$n_1=640,n_2=720,n_3=408$. In each two-coordinate block use
$\log(n_jR)/23$ in the differential basis
$dX/(2Y),X\,dX/(2Y)$. The three fixed collision rows form a
matrix $B$ with
\[
 \bar B=-\begin{pmatrix}
 4&11&0&15&7&4\\22&17&2&8&3&11\\20&4&21&20&12&22
 \end{pmatrix}.
\]
These rows are independent. Exact Jacobian arithmetic also verifies
$D_1=D_2+D_3-D_4$. To compute the rows, translate $n_jD_i$ to
an effective degree-two divisor with distinct integral support roots
and unit ordinates. A root displacement is $-U(x_0)/U'(x_0)$
modulo $23^2$, so first-order tiny integrals give the displayed
six-coordinate logarithms.

For the branch with $z(0)=1$ at every place the moving rows have the
convergent expansion
\[
 B+(h/23)(C_1+hC_2+h^2C_3+\cdots).
\]
The differential pullbacks along the signed-sum branch are
$(-2s/\delta,2r)dc/H_t$. They can be verified on the canonical
intersection of three quadrics: the Jacobian determinant of the
cover equations in $t,r,s$ is $\pm2H_t$. Hence the anti-invariant
canonical forms are spanned by $r\,dc/H_t,s\,dc/H_t$.
Tracing the support of the signed divisor formula at two collision
points with independent $(r,s)/H_t$ fixes the constants and signs.
Here $H_t,r,\delta$ are units on the disks. Thus the unintegrated
series are integral and $nC_n$ has integral coefficients for every
$n\geq1$.

Subtract the fixed rows and remove three factors $h/23$ from the
six-by-six determinant. The reduced determinant has an exact zero
of order at least two at $h=0$. For this assertion, write the three
derivative rows as $u_i(h)R(h)+v_i(h)S(h)$, where the scalars
$u_i,v_i$ are the same rational functions across the embeddings.
One can use $u_i=\sqrt{t_i}/H_t$ and
$v_i=q_2/(\sqrt{t_i}H_t)$, putting the dependence on $z,1/z,\delta$
and $n_j$ in the common embedding vectors $R,S$. A nonzero constant
combination kills both $u_i(0)$ and $v_i(0)$. After integration and
division by $h$, this combination has zero constant term and linear
term in the original two-plane. The other two derivative rows span
that plane at zero. Both the constant and linear determinant terms
therefore vanish. The first nonzero term is
\[
 \det\begin{pmatrix}\bar B\\
       \bar C_1+h\bar C_2+h^2\bar C_3\end{pmatrix}
        =7h^2\pmod{(23,h^3)}.
\]
It controls the full punctured disk. A coefficient of degree $k\geq3$
in this determinant has denominator valuation at most
$\lfloor k/22\rfloor$, since it is built from three integration
indices satisfying $\sum(n_i-1)=k$ and
$v_{23}(n_i)\leq(n_i-1)/22$. At $h\in23\mathbf Z_{23}\setminus\{0\}$
its valuation is strictly greater than $2v_{23}(h)$. The leading
term cannot cancel. This proves independence for equal local signs,
including their common global negative.

\subsubsection{Mixed signs and the two places above 1181}

Up to changing the global sign of $z$, a mixed pattern has one
negative two-coordinate block. Write $S_j$ for its block sign matrix
and $u=h/23$. The determinant of the six rows reduces to
$\det\binom{\bar B}{(\bar B+u\bar C_1)S_j}$. For the three choices
it is
\[
 u(11+20u),\qquad u(21+18u),\qquad u(10+3u)\pmod{23}.
\]
The exact constant term vanishes because at $u=0$ the bottom rows
differ from the top only within one two-dimensional block. Its
quotient by $u$ is an integral convergent series on $\mathbf Z_{23}$,
as each integrated coefficient carries $23^{n-1}/n$. Its constant
is a unit. Thus a dependence with $h\ne0$ requires
$u\equiv19,18,12$ respectively. In each case the six-row matrix
has rank five modulo 23, with normalized left null vector
\[
 (8,7,8,15,4,1),\quad(12,5,22,10,3,1),\quad(19,1,21,13,8,1)
\]
in the order $D_1,D_2,D_3,A_1,A_2,A_3$.

A dependence modulo torsion gives a primitive integer relation, whose
reduction is a nonzero multiple of one of these vectors. The torsion
has trivial 23-primary part because at a split place the good
reduction order is prime to 23 and the formal logarithm is injective.
The relation must therefore vanish in every reduction quotient of
order 23. Use the places over 1181 with $\theta=99,198$. A rational
pole of $c$ is impossible: the three ratios of singleton squares to
$c$ would split the squarefree residue cubic
$144X^3-240X^2-80X+240$, which has only one root there. The
Jacobian reduction orders are $1383036$ and $1338324$, each divisible
exactly once by 23. Multiplication by the order divided by 23 projects
to its order-23 subgroup. Relative to the projection of $D_1$, the
fixed rows have coordinates $(1,11,21)$ and $(1,16,17)$.

At each place enumerate every residue $c$ for which
$c,(25-4c)/6,(c-\theta)/\delta$ are squares and the singleton
cubic splits into three squares, with multiplicity allowed. Enumerate
all ordered signed roots $(a,b,d)$, evaluate the moving divisor
formula, and allow both signs of $z$ independently at each place.
The number of distinct residue tuples $(c,a,b,d)$ satisfying each
possible relation is
\[
\begin{array}{c|rrr}
 &\text{first}&\text{second}&\text{third}\\ \hline
 \theta=99&8&20&16\\
 \theta=198&18&16&22\\
 \text{intersection}&0&0&0
\end{array}.
\]
All the divisors in this finite calculation remain in the nonspecial
Mumford chart. Since $c,a,b,d$ are rational, their residue tuples
must agree at both places. No remaining relation is possible. This
proves that every noncollision beta lift would force six independent
global Prym classes. All divisor, logarithm, determinant and joint
reduction data are checked by \cert{check_type13_beta_prym.m}.

\subsubsection{The unconditional upper bound five}

The quartic factor in the Prym quintic defines an absolute field $L$
with polynomial
\[
\begin{split}
 P(X)={}&X^{12}-4X^{11}-68X^{10}-18698X^9+336998X^8-500324X^7\\
       &+7039932X^6-194015520X^5+1041300636X^4-3497537936X^3\\
       &+20028904184X^2-109223576168X+269574001992.
\end{split}
\]
Its signature is $(2,5)$ and its absolute discriminant is
\[ D_L=172707117717588564754287819295488000000. \]
The unconditional explicit-formula argument for the class group uses
the general-signature formula in \cite[Section~4]{miller}, with
support $X_0=2^{31}$ and
\[
 F(x)=\frac{\max(1-|x|/\ell,0)}{\cosh(x/2)},\qquad\ell=\log X_0.
\]
If the primes of norm at most $b=48000001$ generated a proper
subgroup, a prime-order cyclic quotient would give an everywhere
unramified extension split at all these primes and at the real
places. Subtracting its zeta explicit formula from that of $L$, and
dividing by the degree minus one, gives
\[
 \log D_L\geq A_{\rm arch}+S.
\]
Here each zero contributes nonnegatively throughout the critical
strip. The triangle has nonnegative Fourier transform, and so does
$\cosh(ax)/\cosh(x/2)$ for $|a|\leq1/2$, which proves this without
a hypothesis on zeros. A lower archimedean bound is
\[
 A_{\rm arch}\geq12\left(\gamma+\log(4\pi)
                  -\frac{\pi^2}{4\ell}\right)
                     +2\left(1-\frac{2\log2}{\ell}\right).
\]
A prime-power norm $t=q^m\leq X_0$ has unsigned weight
$4\log(t)(1-\log(t)/\ell)/(m(t+1))$. Its character multiplier
is one for $t\leq b$, and at least $-1$ elsewhere. Counting prime
ideal powers and taking rational interval lower bounds gives
$S>50.333$ and $A_{\rm arch}+S-\log D_L>0.080$, a contradiction.
The computation groups norms in 256 subdivisions of each dyadic
interval. For rational primes above $\sqrt{X_0}$ outside the
polynomial index, only degree-one primes contribute, and their number
is the number of simple polynomial roots. All other primes are
factored in the maximal order.

PARI's exact ideal relation test reduces every prime ideal of norm
at most $b$ to a factor base with norms at most 1000. Only the
outer rational-prime loop is divided among parallel workers. This
proves 1000 as the required generating bound independently of a
provisional class group. The bound is applied only after identifying
the absolute field with $L$. Other fields use their Minkowski bound,
which is 419 for $K$. The ordinary class groups are
$\operatorname{Cl}(L)=(\mathbf Z/2)^2$ and
$\operatorname{Cl}(K)=0$. The Prym two-Selmer group has dimension
six and its rational two-torsion dimension one. Therefore
$\operatorname{rank}J(K)\leq5$, contradicting the six independent
classes forced above. This completes Type 13.
The field and rank checks are \cert{check_type13_beta_prym_rank.m}.
The explicit-formula and ideal-relation computations are
\cert{type13_prym_class/prime_counts.cpp},
\cert{type13_prym_class/check_bound.py} and
\cert{type13_prym_class/cover.py}.

\subsection{Type 9 and a height comparison within a split fiber}\label{app:type9-details}

The singleton-root cubic in Appendix~\ref{app:type9} has a projective
linear identification with
\[
 E:y^2+y=x^3+x^2+295x+1109.
\]
For coordinates $[R:B:W]$ on that cubic and standard projective
coordinates on $E$, forward and inverse matrices are
\[
 L=\begin{pmatrix}-44&-84&347\\-1088&3102&4064\\16&6&-28\end{pmatrix},
 \qquad J=\begin{pmatrix}2884&7&36757\\-896&112&5152\\1456&28&5908\end{pmatrix}.
\]
Their product is a nonzero scalar matrix. The full group is
$E(\mathbf Q)=\mathbf ZG$ with $G=(31,202)$ and trivial torsion,
certified with unconditional class-group bounds by
\cert{check_type9_group.m}.

Three distinct rational roots at fixed finite $b=B/W$ are three
collinear points of the plane cubic. Their images are collinear on
the Weierstrass model, so they have indices $n_1+n_2+n_3=0$.
If $M=\max|n_i|$, one of the other absolute indices is at most
$M/2$. The point at infinity on $E$ maps to $[7:112:28]$, so the
fiber containing it is exactly $b=4$. We treat it separately.

Write $b=N/D$ primitively with $D>0$, $H=\max(|N|,D)$, and let
$H_C$ denote the maximum primitive coordinate of a point on the
original cubic. Its integral equation for $Dr$ has leading
coefficient 64, so $64Dr$ is an integer. After scaling by $H$, the
remaining coefficients are bounded by $1040/64,4396/64,4707/64$.
The root bound gives $|Dr|<80H$. The coordinates
$[64Dr:64N:64D]$ and their ratio imply
\[
 H\leq H_C\leq5120H.
\]
The maximum absolute row sums of $L,J$ are less than 9000 and
40000. Thus the primitive projective height $H_E$ on $E$ satisfies
$H_E\leq9000H_C$ and $H_C\leq40000H_E$.

For a finite elliptic point write $x=a/t^2$, $y=c/t^3$ in reduced
form, $t>0$, and $U=\max(\sqrt{|a|},t)$. The primitive projective
coordinates are $[at:c:t^3]$. The equation bounds
$|c|^2\leq |c|U^3+1406U^6$, giving the upper estimate below.
For the lower estimate, $|x|\leq100$ gives $U^3\leq1000t^3$.
The range $x<-100$ is impossible because the cubic is less than
$-1/4$, while $x>100$ implies $|y|>x^{3/2}/2$. Hence
\[
 U^3/1000\leq H_E\leq40U^3.
\]
Put $h(R)=\log_2\max(|a|,t^2)$. For any two points in one fiber,
these estimates show that
\[
 2^{3(h(R_i)-h(R_j))/2}
   \leq1000\cdot40\cdot9000\cdot40000\cdot5120<2^{60}.
\]
Their naive heights differ by less than 40, uniformly in $b$.

For the doubled canonical height
$\lambda(R)=\lim4^{-k}h(2^kR)$, the explicit duplication forms are
\[
\begin{aligned}
 F&=A^4-590A^2B^2-8874AB^3+82588B^4,\\
 D_2&=4A^3B+4A^2B^2+1180AB^3+4437B^4.
\end{aligned}
\]
There are integral Bézout identities
$dV^7=S_VF+T_VD_2$ for $V=A,B$, with
$d=2079016875$ and the sum of absolute coefficients of each pair
$(S_V,T_V)$ less than $2^{46}$. They imply
$\gcd(F,D_2)\mid d$ for primitive $A,B$. Taking the larger of
$|A|,|B|$ bounds the unreduced maximum below by
$d\max(|A|,|B|)^4/2^{46}$. The coefficient sums of $F,D_2$ are
less than $2^{17}$. Removing the gcd and telescoping yields
\[
 -46/3\leq\lambda(R)-h(R)\leq17/3.
\]
The numerator or denominator of $x(32G)$ has bit length 149, so
$\lambda(G)\geq(148-46/3)/32^2>1/8$.
By canonical-height quadraticity~\cite[Chapter~VIII, \S9]{Silverman},
heights of points within the fiber differ by
less than $40+(17+46)/3=61$. Comparing indices of maximum and
minimum absolute value gives
$3M^2\lambda(G)/4<61$. Since $3\cdot26^2/32>61$, all three
indices are bounded by 25. This is an all-height bound on any
split fiber.

Exact projection of the 51 multiples from $-25G$ to $25G$ gives
only two finite parameters occurring more than once,
\[
\begin{array}{c|l}
 b& (n,R/W)\\ \hline
0&(-2,25/4),(1,9/4)\\
4&(-3,49/4),(0,1/4),(3,81/4).
\end{array}
\]
At zero the cubic is $(4R-9)^2(4R-25)$, while at four it is
$(4R-1)(4R-49)(4R-81)$. These are also the only repeated-root
possibilities. Indeed the remaining factor of the discriminant has
unit leading coefficient at 7 and nonzero residues
$2,1,6,1,3,2,6$ on $\mathbf F_7$. A rational root would be
7-integral and reduce to a root. The only distinct split fiber is
therefore $b=4$, where the doubleton centered at 2 shares $3/2$
with the central chain. The Bézout forms, heights, exact multiples
and discriminant test are \cert{check_type9_exclusion.py}.

The ordered-root curve is the complete intersection in
$[r:s:B:W]$ defined by $f(r,B,W)=0$ and
$(f(r,B,W)-f(s,B,W))/(r-s)=0$. It is smooth of genus four.
At a double root the base derivative of $f$ and the derivative
$f''/2$ of the divided difference supply independent tangent
conditions. Above infinity the discriminant is $-675/4$, so there
are no rational points there. The two finite fibers above give all
nine rational points, including the three orderings at the repeated
fiber. This adds a complete rational-point consequence to the
height argument.

\subsection{Type 3 and a complete rational-image calculation}\label{app:type3-details}

In Appendix~\ref{app:type3}, the last center constraint has
$S_0=(20b-8c+27)/6$. If $S\ne S_0$, taking the divided residuals,
their resultant in $S$, and its resultant with the conic in $b$
leaves the factors
\[
\begin{split}
 &(4c-25)(4c-9)(4c-1)^4(80c^2-280c+281)\\
 &\hspace{12mm}\cdot(144c^2-616c+3025)(144c^2-328c+4489).
\end{split}
\]
The three quadratics have negative discriminants. The linear cases
are collisions or the nonsquare values $b=15/32,77/32$. Thus the
rational parameterization with $S=S_0$ covers every valid input.
For its parameter $z>0$, use the polynomials $L,A,B,C,V,D$ of
Appendix~\ref{app:type3}. If $z=N/M$ primitively with $M>0$, the
squarefree parts $d,k$ of $L(N,M)$ and $NM$ satisfy
$d\mid690$, $k\mid138$. The complete conditions are that
\[
 dL,\ dA,\ dB,\ dC,\ kNM,\ 3kV,\ D
\]
are integer squares. The resultant of $L$ and $A$ has prime support
$\{2,3,5,23\}$, and the outer coefficients of $V$ have support
$\{2,3,23\}$, which proves those finite squareclass ranges.
Both projective patches are included in the local test. Modulo
$2^{13}$ the remaining pairs are
$(1,6),(3,46),(115,46),(345,6)$. Imposing $3^7$ leaves just
$(1,6),(115,46)$.

For the first pair $6z$ and $L(z)$ are squares, giving a nonzero
point on
$E_6:Y^2=X(X^2-342X+59616)$. Its rank-zero group has only the
finite point $(0,0)$, contradicting $z>0$.
For the second pair, put $t=1656/z$. It is a nonzero square, and
\[
 10(t^2-57t+1656)\in\mathbf Q^{\times2},\qquad
 15(t^2-131t/8+3381)\in\mathbf Q^{\times2}.
\]
On $E_{10}:Y^2=X(X^2-570X+165600)$ there is therefore a point
with $X=10t$. The rational group has rank one and torsion generated
by $(0,0)$. The point $(160,4000)$ has two-isogeny class 10,
independent of the torsion class 46. Thus the abscissa class 10
forces its difference from this point to be twice a rational point.

Over $K=\mathbf Q(\sqrt{-15})$, put
$r=(57+15\sqrt{-15})/2$. The abscissa character at $(10r,0)$
implies $(t-r)/(16-r)\in K^{\times2}$. One can see the character
homomorphism directly: for three collinear points the product of
their $X-10r$ factors is the square of the line ordinate at $10r$.
Let $q(t)=t^2-131t/8+3381$ and $l=64(16-r)3/5$.
The two square conditions now give a point with $X=l(t-r)$ on
\[
 E:Y^2=X\bigl(X^2+lq'(r)X+l^2q(r)\bigr).
\]
Writing $s=\sqrt{-15}$, its two coefficients are
$45300-18900s$ and $488160000s-4186080000$.
The two-Selmer group has order four and torsion is generated by
$(0,0)$. The point $P=(X_0,360X_0)$,
$X_0=7200s-26400$, is non-torsion, so the rank is one.

The subgroup generated by $P$ and torsion has odd index. At the
good place $(19,s-17)$, both $X(P)$ and $X(P+(0,0))$ reduce to
nonsquares, so neither translate is a double. Elliptic Chabauty for
the map $t=r+X/l$ returns three subgroup points with rational
image, and an index obstruction equal to a power of two (4 in the
stated run). The obstruction is coprime to the subgroup index, so
the rational images are exhausted. They are $\infty$ and 16.
This uses the rational-image index conclusion of elliptic Chabauty,
so an unproved full Mordell--Weil basis is unnecessary.
Every rank and class-group computation uses unconditional proof
bounds. At the only finite image $t=16$, one has
$z=207/2$ and $(a,b,c)=(81/4,25/4,25/4)$. The singleton at
$5/2$ meets the center of the tripleton on the other side.
The elimination and squareclass tables are
\cert{check_type3_exclusion.py}, and the complete rational-image
certificate is \cert{check_type3_exclusion.m}.

\subsection{Type 17 and the primitive dyadic direction}\label{app:type17-details}

The product elliptic curve $E_0$ of Appendix~\ref{app:type17} has
full rational group generated by $G=(11025/16,165375/64)$ and no
torsion. On odd multiples, the cubic-root abscissa character fixes
$d=1/16-r$, where
\[
 w^3+w^2-240w+1120=0,\qquad
 r=(2w^2+5w+140)/3675.
\]
The associated curve $Y^2=X(X+5d)(X-rd)$ has full two-torsion,
Selmer order eight and rank one, with non-torsion point
$P=(d/16,9d^2/16)$. At the two good places above 17 with
$w=6,11$, none of the four translates of $P$ by two-torsion is a
double. The generated subgroup consequently has odd index.
Elliptic Chabauty with index obstruction 16 gives all rational
images $X/d$, namely $\infty,0,-5,1/16$. Only the last is finite
and positive, giving $b=16$, a collision. This exhausts the entire
odd-multiple class.

For the even class, write the product point as $2Q$, put $a=11025$
and $T=X(Q)$. Retain the cubics and quartics $g,A,B$ of
Appendix~\ref{app:type17}, so
\[
 x=\frac{A(T)}{4ag(T)},\qquad x+5=\frac{B(T)}{4ag(T)}.
\]
All three of $g(T),A(T),B(T)$ are nonzero squares. The odd-degree two-cover descents~\cite{BruinStoll} on $y^2=gA$ and $y^2=gB$ have sizes two
and one respectively, each containing the identity. The nonidentity
class in the first set projects to a nonsquare in the cubic field
of $g$. A representative in its cubic-times-quartic étale algebra
is divided by its total norm before projection. For an actual
abscissa class that norm is square. Multiplication of the original
representative by a rational scalar changes the normalized class
by a sixth power, which is square. Thus the normalization is
well defined. The two descents share the cubic-root character
of $T-r$. The second descent makes it trivial, selecting the
identity in the first as well. This proves that $T-\theta$ is
square for a root $\theta$ of $A$, and that $Q$ is an even
multiple of $G$.

On the integral equation of $E_0$, the generator has formal
parameter $-4/15$. The formal subgroup on $4\mathbf Z_2$ is
closed under addition and duplication starts with twice the
parameter. Therefore for $s=-T/Y(Q)$,
\[
 k=v_2(s)\geq3,\qquad v_2(T)=-2k.
\]
Take $K=\mathbf Q(w)$, $w^4-194w^2-96w+8302=0$, and
$\theta=3(2w^3+4w^2-240w-532)$. The monic cubic
$h(U)=A(U)/(U-\theta)$ is integral and irreducible over $K$.
The two square conditions give a global point $R=(T,Y')$ on
$E:Y'^2=h(T)$.

Its two-division root field has absolute polynomial
\[
 Z^{12}-472Z^{10}+75282Z^8-4450612Z^6+
                  64570537Z^4-277288596Z^2+234027000.
\]
It is totally real, with discriminant
$58160039859394517078535559761100800000$. The unconditional
explicit-formula estimate used above, now with support and
norm bound $2^{20}$, has margin greater than $0.262$.
Every prime-power term in the support has multiplier one in the
hypothetical cyclic quotient. PARI's exact ideal relations reduce
the generating set to norm 1000. After field identification this
gives class group $\mathbf Z/2$ and
$\dim_{\mathbf F_2}\operatorname{Sel}_2(E/K)=2$.

There is one dyadic place of $K$, totally ramified of degree four.
Write $D=K_v$, let $\pi$ be a uniformizer and $\mathfrak p$ its
maximal ideal. The integral basis is $1,\pi,\pi^2,\pi^3$.
The cubic $h$ has exactly one root in $D$, so
$\dim E(D)/2E(D)=5$, four from the local Lie group and one from
local two-torsion. Put $m=105\cdot2^{12}$ and
$\Lambda=m\log E(D)$. Five explicit points with independent local
Kummer images span this quotient. Their scaled logarithms span
$\Lambda/2\Lambda$, since a torsion point has logarithm zero.
For each point, the formal parameter of its $m$-multiple has
$v_\pi\geq44$. On the displayed integral model the differential
has integral even coefficients, so its logarithm is
$z+O(z^3)$ with dyadically integral coefficients. Its leading binary
coordinates may therefore be obtained from the formal parameter.
The five rows in the basis
\[
 (2^{12},2^{11}\pi,2^{11}\pi^2,2^{11}\pi^3)
\]
are integral with residue rank four. It follows that
$\Lambda=2^{11}\mathfrak p$. Expressing the two localized global
Selmer classes in the five local Kummer generators and then using
their logarithm residues gives a two-dimensional image in
$\Lambda/2\Lambda$. Adjoining the rational axis $(1,0,0,0)$
raises the rank to three. This simultaneously proves injectivity
of the Selmer-to-logarithm map and exclusion of that primitive
axis.

The chain equations force exactly the excluded direction. For
$z=-T/Y'$ the ratio satisfies
\[
 (z/s)^2=\frac{g(T)}{h(T)}\in1+2^{2k}\mathcal O_D.
\]
Both cubics are monic integral and $v_2(T)=-2k$, which proves this
congruence by dividing them by $T^3$. Of $z/s-1,z/s+1$, the one
with smaller valuation has valuation one, using the normalization
$v_2(2)=1$. Thus for one sign
\[
 z=\pm s(1+2^{2k-1}\eta),\qquad\eta\in\mathcal O_D.
\]
The same leading congruence holds for $\log_E R$. Consequently
$2^{-k}m\log_E R$ is primitive in $\Lambda$ and its residue is
$(1,0,0,0)$.

To apply the global Selmer obstruction, divide $R$ by the largest
possible global power of two, writing $R=2^jR_0$. The nonzero
logarithm and finite generation make $j$ finite. A globally
primitive point has nonzero Kummer class, hence a nonzero residue
in $\Lambda/2\Lambda$ by injectivity of the map just computed.
Therefore $m\log_E R_0$ is primitive, forcing $j=k$.
Its residue would be the excluded axis. This contradiction
exhausts the even branch and proves Type 17 at all rational
heights. The two-cover projections, local generator independence,
logarithm valuations and final three-dimensional rank assertion
are all checked by \cert{check_type17_even_branch.m}.
The odd cover is \cert{check_type17_odd_branch.m}, and the
unconditional relation range is
\cert{check_type17_factor_base.gp}, together with the
\texttt{--type17} mode of
\cert{type13_prym_class/check_bound.py}.
